% =====================================================================
% S6 background-prose draft — Derived Consequences of the Structured Model
% Working structure: IWAI 2026, Hallgren & Hyland, 2026-05-29
% Author of draft: Claude sub-agent (one of 7 parallel §-drafts)
% Target length: 500–800 words
% Clusters drawn on: C4 social_network_epistemology (derived-results subset);
%                    C5 echo_chambers (Nguyen bubble-vs-chamber distinction)
% Overlap note: §3 author handles the *primitives* (Bala–Goyal, DeGroot, MAB).
% This draft handles the *derived consequences* — Zollman tradeoff, Kitcher–
% Strevens DoCL + paradoxical stubbornness, Nguyen bubble vs chamber.
% --------------------------------------------------------------------
% TODO hand-rename anchors before merging the cluster bibs into refs.bib:
%   - Zollman 2007 (network epistemology / less-is-more) — anchor MISSING from
%     cluster_C4; rec key zollman2007NetworkEpistemology;
%     and Zollman 2010 (transient diversity) — anchor MISSING; rec key
%     zollman2010TransientDiversity. Both cited below as
%     \citep{zollman2007NetworkEpistemology, zollman2010TransientDiversity};
%     need to be added to refs.bib (not yet present).
%   - Kitcher 1990 "The Division of Cognitive Labor" — anchor MISSING from
%     cluster_C4 (referenced in-text by Romero2017_2); rec key
%     kitcher1990DivisionCognitiveLabor.
%   - Strevens 2003 "The Role of the Priority Rule in Science" — anchor
%     MISSING; rec key strevens2003PriorityRule.
%   - Muldoon & Weisberg 2011 "Robustness and Idealization in Models of
%     Cognitive Labor" — anchor MISSING from cluster_C4 (cited inside
%     Romero2017_2's references); rec key muldoonWeisberg2011Robustness.
%   - Nguyen 2018 (Episteme, cluster key Nguyen2018_1) — rec key
%     nguyen2020EchoChambers (the 2020 reprint is the canonical citation).
%   - Longino 1990 "Science as Social Knowledge" (Longino1990_10 in cluster) —
%     rec key longino1990ScienceSocialKnowledge; this is the CCE / four-norm
%     anchor.
%   - Romero 2017 (cluster key Romero2017_2) — rec key
%     romero2017NoveltyReplicability; needed for the priority-rule /
%     replicability vice.
%   - Pollock 2024 (cluster key Pollock2024_3) and Turner 2023 (cluster key
%     Turner2023_2) follow Nguyen's distinction; cite from C5 as written.
% Style reference: paper/sections/02_related_work.tex (Bourdieu-style
% \paragraph blocks, \citep / \citet, our internal macros).
% =====================================================================

\section{Derived Consequences of the Structured Model}
\label{sec:derived-results}

The structural extension developed in Section~\ref{sec:experiments} is
not introduced for its own sake. It earns its place because three
otherwise-disparate results in network epistemology --- the Zollman
tradeoff between consensus speed and asymptotic accuracy, the
paradoxical-stubbornness finding from the division-of-cognitive-labor
literature, and Nguyen's distinction between epistemic bubbles and
echo chambers --- recover as derived consequences of one mechanism: a
trust-precision network propagating candidate \emph{structural}
revisions rather than scalar opinions. We position each result in turn
and indicate which observable in our simulations it indexes.

\paragraph{The Zollman tradeoff and the value of inefficiency.}
\citet{zollman2007NetworkEpistemology, zollman2010TransientDiversity}
showed, in a multi-armed-bandit network of Bayesian scientists, that
\emph{denser} communication graphs converge faster but more often
lock onto the worse arm; sparser graphs sustain transient diversity
long enough for the inferior hypothesis to be discarded, and so reach
higher asymptotic accuracy. The mechanism is now well understood as a
spectral one: the rate at which beliefs synchronise on a graph is set
by its algebraic connectivity (the Fiedler value $\lambda_2$ of the
graph Laplacian)
\citep{Griparic2021_7, Montijano2017_11, Tan2008_14}, so larger
$\lambda_2$ compresses the transient phase during which dissent can
falsify a leading hypothesis. The structured model inherits this
argument unchanged at the topological level, but reroutes its
substantive content: the channels carry \emph{precision-weighted
candidate revisions} of the generative model, not scalar credences,
so the speed–accuracy tradeoff applies to structural revision events
themselves. This vindicates the Feyerabendian-pluralist intuition
\citep{Müürsepp2019_7} in mechanistic form --- a measured inefficiency
in the trust-precision network is the substrate condition for the
critical-discursive uptake that
\citet{longino1990ScienceSocialKnowledge} requires of an epistemically
healthy community --- and supplies the formal claim Longino's
critical contextual empiricism has been criticised for leaving
underspecified \citep{Novaes2025_1, Ajdari2023_6, Borgerson2011_5}.

\paragraph{Division of cognitive labour and paradoxical stubbornness.}
\citet{kitcher1990DivisionCognitiveLabor} and
\citet{strevens2003PriorityRule} argued that the priority rule
incentivises credit-driven scientists to spread across research
programmes, producing a community-level distribution of effort closer
to the social optimum than truth-driven scientists would individually
choose. \citet{muldoonWeisberg2011Robustness} relaxed the
god's-eye information assumption: in their local-information variant,
agents that appear \emph{paradoxically stubborn} --- unwilling to
abandon a low-success-probability programme --- in fact regularise
the system, preventing premature convergence on whichever programme
happened to lead first. \citet{romero2017NoveltyReplicability}
qualifies the optimistic Kitcher--Strevens reading by noting that the
priority rule rewards novelty regardless of replicability and so
\emph{disincentivises} the very direct-replication work that would
discipline structural revisions. In the structured model these
findings appear as a single statement: an agent's resistance to a
structural revision is proportional to the precision of its current
generative model, so agents with strong structural utility on their
incumbent paradigm act as protective elements against premature
structural collapse. They are not uniquely irrational; their
structural commitments simply require correspondingly stronger
cross-channel evidence to revise. The Muldoon--Weisberg paradox and
the Romero replication-vice fall out as two readings of the same
asymmetry --- one casting it as virtue (transient diversity
preserved), the other as vice (replication suppressed) --- and the
model exposes the parameter that interpolates between them.

\paragraph{Bubbles versus chambers: failure modes of cross-agent channels.}
\citet{nguyen2020EchoChambers} draws a distinction that, in our
framework, is structural rather than rhetorical. An
\emph{epistemic bubble} is a failure of \emph{trust topology}:
relevant voices are omitted, so the cross-agent channel is sparse or
disconnected, and the bubble pops on simple exposure. An
\emph{echo chamber} is a failure of \emph{trust precision}: the
discrediting beliefs assigned to outsiders pre-emptively downweight
incoming cross-surprisals, so the chamber survives exposure and can
even use it for reinforcement
\citep{Pollock2024_3, Turner2023_2, Clinkenbeard2025_4}. In the
scalar-opinion setting this distinction must be argued for
philosophically; in the structured-belief setting it is recovered as
the formal separation between sparsity of $\trustmat$ and the
precision weights it carries on its non-zero entries. Crucially,
because the channels now carry candidate \emph{structural} revisions
rather than scalar opinions, the bubble--chamber distinction becomes
more honest: a bubble withholds candidate revisions; a chamber
admits them and then assigns them zero precision. The two failure
modes call for different repairs --- coverage versus
precision-recalibration --- and the model makes the repair
prescription mechanistic rather than virtue-theoretic.
